\subsection{Approximate Line Search}

\begin{frame}{Approximate Line Search: Explicación Gráfica}
    \begin{block}{Concepto principal}
        [Aguilar: Explica las condiciones de Wolfe o Armijo y cómo visualmente buscan un paso que no sea ni muy corto ni muy largo].
    \end{block}
    \begin{itemize}
        \item {[Idea gráfica 1]}
        \item {[Idea gráfica 2]}
    \end{itemize}
\end{frame}

\begin{frame}{Approximate Line Search: Pseudocódigo}
    \begin{algorithmic}[1]
        \State \textbf{Input:} [Entradas y parámetros $c$ y $\rho$]
        \State [Lógica de Backtracking]
        \State \textbf{Return} [Tamaño de paso $\alpha$]
    \end{algorithmic}
\end{frame}

\begin{frame}{Approximate Line Search: Ejemplo Paso a Paso}
    \begin{exampleblock}{Planteamiento del problema}
        [Aguilar: Función de prueba multivariable y dirección de descenso elegida].
    \end{exampleblock}
    \begin{itemize}
        \item \textbf{Prueba $\alpha=1$:} [¿Cumple la condición de Armijo?]
        \item \textbf{Backtracking:} [Reducción del paso]
    \end{itemize}
\end{frame}

\begin{frame}[fragile]{Approximate Line Search: Código}
\begin{lstlisting}[language=Python]
# [Aguilar: Reemplaza esto con tu código real]
def approx_line_search_ejemplo():
    pass
\end{lstlisting}
\end{frame}

\begin{frame}{Approximate Line Search: Análisis de Convergencia}
    \begin{alertblock}{Eficiencia}
        [Aguilar: Menciona por qué no buscar el mínimo exacto acelera el proceso global (ahorro de evaluaciones)].
    \end{alertblock}
    \begin{itemize}
        \item {[Garantías de convergencia global (Teorema de Zoutendijk)]}.
    \end{itemize}
\end{frame}